% P03: INSERTION - new subsection + figure, placed immediately BEFORE
% \section{Capabilities and Remaining Limits}\label{sec:sns_discussion}
\subsection{Perturbation Robustness and Pathway Deletions}\label{sec:w2l_perturbation}

I tested the coupled W2L walker the way the perturbation literature tests neuromechanical models: the tuned conductances stay frozen, the perturbation is applied to the body or to the network, and recovery is judged from the gait itself (Figure~\ref{fig:w2l_perturbation}). The frozen configuration steps with near-equal stride counts on the two legs (17 left, 16 right in the 16-second window). Under segment-mass changes --- thigh and distal-leg masses each scaled by $+10$, $-10$, $+20$, and $-20$ percent, and the whole leg by $\pm 10$ percent --- stepping persisted in every condition (33--38 strides per leg over the longer test window, pelvis never below \qty{0.90}{\meter}), although in nine of ten conditions the pelvis lean exceeded the \ang{45} fall criterion of the evaluation gate, so the mass robustness of the posture, as distinct from the stepping, remains open. An inhibitory tonic stimulus applied to the rhythm-generator layer for one second left the stride counts essentially unchanged (18/17 mild, 19/17 strong), and a \qty{40}{\newton} fore-aft push at the pelvis likewise (18/18 and 18/16). Removing individual pathway classes then located the load-bearing connections: deleting the crossed heel-to-contralateral-flexor edge, the Ia reciprocal chain, the Renshaw cells, the heel-contact cascade, or the commissural coupling left stepping at or near the baseline counts (16/16, 13/13, 16/14, 17/15, and 13/12 respectively), whereas removing the Ib load feedback stopped stepping entirely (0/0). The same configuration walks with the harness weight support re-tuned down a ladder from 70 to 10 percent of body weight, with a fresh-process verification at every step. These tests concern supported progression under the harness; unsupported walking is not established.

\begin{figure}[p]
\centering
\includegraphics[width=\textwidth,height=0.52\textheight,keepaspectratio]{w2l_perturbation_panels.pdf}
\caption[W2L perturbation and pathway-deletion panels]{Perturbation tests of the corrected W2L walker in MuJoCo with all conductances frozen (harness at \qty{25}{\percent} weight support). (A--C) Left and right knee flexion (solid and dashed; heel contact shaded). (B, C) A one-second inhibitory tonic stimulus to the rhythm-generator layer (yellow band), mild and strong; stepping persists through and after the window. (D, E) Fore and aft \qty{40}{\newton} pelvis pushes at $t=\qty{8}{\second}$ (dotted line): fore-aft center-of-mass position (solid) and pelvis tilt (dashed). (F) Strides per leg when one pathway class is removed: most removals are compensated, while commissural coupling and Ib load feedback are load-bearing.}
\label{fig:w2l_perturbation}
\end{figure}
